\chapter{Marco Teórico}
Debe mencionar definiciones, modelos geológicos, discusiones u otros aspectos que no hayan quedado explicados en la introducción.

% ---------------------------------------------------------------------------
% EJEMPLOS DE ORGANIZACION DEL TEXTO
% Los comandos con * crean titulos sin numero y no los agregan al indice.
% Quitar este bloque cuando ya no sea necesario como guia.
% ---------------------------------------------------------------------------
\section{Ejemplo de organización del texto}\label{sec:ejemplo-organizacion}

Las secciones creadas con \verb|\section{...}| se numeran y aparecen en el
índice general. También pueden llevar una etiqueta para citarlas desde otro
lugar, por ejemplo: véase la sección~\ref{sec:ejemplo-organizacion}.

\subsection{Subsección numerada}

Una subsección creada con \verb|\subsection{...}| se numera y aparece en el
índice debajo de la sección correspondiente.

\subsection*{Subsección sin número}

Al agregar un asterisco, como en \verb|\subsection*{...}|, el título no lleva
número y no aparece automáticamente en el índice. Esto resulta útil para
aclaraciones o apartados breves.
